\section{Introduction}
\label{sec:intro}

Claims organisations decide every day which documents to request from a claimant. The decision is governed by statutes, policy terms and internal guidelines whose conditions refer to the facts of the case, and a single fact can change it. If a stolen watch was separately insured above a threshold, the insurer needs proof of its value; if it was not, a receipt is unnecessary; if nobody knows yet, asking for the receipt is premature, and the right step is to ask whether the watch was insured. Errors in either direction are costly: unnecessary requests delay claims and burden claimants, while missed requests let a claim proceed without the evidence that a payment or refusal requires \cite{eling2022impact}.

Language-model agents are increasingly asked to make such decisions. The prevailing design places the organisation's rules in the agent's prompt and lets the model decide \cite{yao2025taubench,xiao2024flowbench}. The model then judges the case and applies the rules in a single generation, so neither decision is recorded separately, and repeated runs can disagree \cite{atil2025nondeterminism,blairstanek2025unstable}. The alternative is to let code apply the rules. Rule engines and logic programs do this traceably \cite{sergot1986british,merigoux2021catala,omg2024dmn}, and neuro-symbolic systems let a model translate a problem for a solver \cite{pan2023logic,olausson2023linc}. Both, however, assume that every relevant fact is known or treat unknown facts as false, whereas at intake many facts are not yet known.

This article asks one question: \emph{should a language-model agent apply an organisation's evidence rules itself, or should code apply them to the agent's verified judgments of the case?} The question matters for three reasons. Organisations must decide where to place this responsibility when they deploy agents in regulated processes \cite{dumas2023aiaugmented,hevner2023research}. Executing rules promises guarantees that an agent's own judgment cannot give: no premature requests, changes only where facts change, and a checkable reason for every request. And it is not known whether those guarantees translate into better requests, or what they cost. An earlier version of this work reported fewer wrong checklist changes for executed rules, but reviewers rightly noted that it used one model, single runs and home-made baselines, and that some of its comparisons favoured execution by construction.

We answer the question with \casepath, an agent that separates the two responsibilities (Figure~\ref{fig:overview}). A language model reads the claim and records, for every rule condition, whether it is true, false or unresolved, quoting the text that supports the verdict; a second model call verifies these judgments, and code admits only verdicts whose quotations occur verbatim in the sources. Code then executes the organisation's process rules: it requests the documents that are certainly required, holds those whose requirement depends on an unresolved fact behind a clarification question, and skips the rest. We show that such requests are never premature, identify a case in which the three-valued logic of \casepath's process engine holds documents that every answer would require, and correct it by evaluating requirements over all completions of the open facts.

\begin{figure}[t]
\centering
%% Figure 1: the running example through CasePath's four stages (TikZ, editable).
\begin{tikzpicture}[
  font=\footnotesize,
  box/.style={draw=black!35, rounded corners=2pt, align=left, inner sep=4pt, fill=white},
  model/.style={box, draw=blue!55!black},
  code/.style={box, draw=black!70},
  lab/.style={font=\scriptsize\sffamily, text=black!60},
  arr/.style={-{Stealth[length=4pt]}, black!60, line width=0.6pt}]

\node[box, text width=3.05cm] (claim) {\textbf{Claim}\\[1pt]
``\ldots a laptop and a watch were stolen. The watch may be separately scheduled in the policy, but I have not found the policy schedule yet.''};

\node[model, text width=3.55cm, right=0.55cm of claim.north east, anchor=north west] (interp) {\textbf{1\enspace Interpret}\hfill{\scriptsize model}\\[1pt]
\emph{watch separately insured} $=\U$\\
quote: ``may be separately scheduled''};

\node[model, text width=3.55cm, below=0.25cm of interp] (verify) {\textbf{2\enspace Verify}\hfill{\scriptsize model}\\[1pt]
second call accepts or rejects each judgment};

\node[code, text width=3.55cm, below=0.25cm of verify] (admit) {\textbf{3\enspace Admit}\hfill{\scriptsize code}\\[1pt]
keeps a verdict only if its quote occurs verbatim and it was accepted; otherwise $\U$};

\node[box, text width=3.75cm, right=0.55cm of interp.north east, anchor=north west] (rule) {\textbf{Rule} \hfill{\scriptsize from policy terms}\\[1pt]
if the watch is separately insured above CHF 1{,}000: establish its value by a receipt \emph{or} an expert valuation};

\node[code, text width=3.75cm, below=0.25cm of rule] (exec) {\textbf{4\enspace Execute}\hfill{\scriptsize code}\\[1pt]
$\T$: \textbf{request} the receipt (fewest documents)\\
$\F$: \textbf{skip} both documents\\
$\U$: \textbf{hold} both and ask:\\ \emph{``Was the watch separately insured?''}};

\draw[arr] (claim.east |- interp.west) -- (interp.west);
\draw[arr] (interp) -- (verify);
\draw[arr] (verify) -- (admit);
\draw[arr] (admit.east) -- ++(0.25,0) |- (exec.west);
\draw[arr] (rule) -- (exec);
\node[lab, below=0.05cm of exec.south, anchor=north] {every request keeps its trace: source passage $\to$ rule $\to$ quote};
\end{tikzpicture}

\caption{\casepath on the running example. A model judges each rule condition from the claim and quotes its evidence; a second call verifies the judgment; code admits only verbatim quotations and then executes the source-grounded rule, which requests, skips or holds the documents depending on whether the condition is true, false or unresolved.}
\label{fig:overview}
\end{figure}

We then compare \casepath with the main alternatives from the literature---policy text in the prompt, the rules in the prompt, majority voting over runs, stepwise rule application by the model, and two-valued rule execution---under identical inputs. The evaluation uses 90 tenancy claims with complete references, 27 pairs of household-theft cases that differ in one fact, and 18 new pairs in which that fact is left open; every design runs three times with each of two current models, and the analysis plan was fixed before the held-out runs. The answer depends on the model. With \texttt{gpt-6-luna}, the cost-efficient model configured in the \casepath product, executing the rules made about four times fewer unsafe errors on complete request lists (\tenLunaCPxUnsafe{} versus \tenLunaBoneUnsafe{} per claim) and four times fewer wrong checklist changes (\thLunaCPWrong{} versus \thLunaRWrong{}) than letting the model apply them, and it left far more requests unchanged when a fact was open; the original study's model showed the same pattern on changed facts in its single run. With the frontier model \texttt{gpt-6.1-terra}, the model applying the rules in context was as safe as execution on complete request lists and, though not significantly, more accurate on automatically extracted rules, whose errors execution reproduces exactly. Across the three models, \casepath made between 0.21 and 0.56 wrong changes per pair, against 0.27 to 1.52 for the model applying the rules in context. Execution's main cost is deferral: \casepath held 20--30\% of required documents behind questions, most of them because the claims leave one routing fact unstated; reading those open facts as false or as true instead would have produced silent misses or unnecessary requests.

The article makes four contributions:
\begin{enumerate}[leftmargin=1.5em]
\item \emph{A precise formulation} of evidence-request decisions under incomplete case facts as three-valued execution of a source-grounded process, with soundness and locality properties, an incompleteness result for strong Kleene evaluation, and exact deferral by supervaluation, which removes \tenTerraExactCut--\tenLunaExactCut\% of the production engine's deferrals without adding unnecessary requests (Section~\ref{sec:method}).
\item \emph{The \casepath agent}, which combines typed interpretation, independent verification, exact-citation admission and rule execution, released with its experiment harness \cite{casepath2026repo}.
\item \emph{A systematic comparison} with the alternatives in the literature on two datasets, two current models, three runs per case and new unresolved-fact variants, with independent reference checks and exact error-propagation analysis (Sections~\ref{sec:evaluation} and~\ref{sec:results}).
\item \emph{Design knowledge} on when executing rules pays: it makes request quality robust to the model, safe under uncertainty and auditable, but it is only as accurate as its rules and its treatment of unstated facts (Section~\ref{sec:discussion}).
\end{enumerate}
Section~\ref{sec:related} reviews related work, Section~\ref{sec:method} presents \casepath, Section~\ref{sec:evaluation} describes the evaluation and Section~\ref{sec:results} its results; Section~\ref{sec:discussion} discusses implications and limitations.
